\section{Integration of Assembly Wrapping}\label{sec:wrapping}
\subsection{Existing Technique and Implementation Scope}
Assembly wrapping is an existing technique based on alternative
instruction interpretations~\cite{assembly_wrapping}. The contribution
here is its integration into Zone and empirical assessment. The
instruction view obtained by static decoding can differ from the
stream followed during execution; representation measurements must
state which view they use.

The implementation now supports both x64 and x86. The x86 path was
added after the first benchmark snapshot and is assessed against a
separately identified version. Register width, instruction encodings,
and addressing require architecture-specific treatment; an x64 outcome
is not transferred to x86 by assumption.

\subsection{Integration and Assessment Boundary}
Wrapping depends on an adequate instruction interpretation, data
preparation, and consistent handling of supported references. Some
regions may remain outside the main wrapped representation, and
construction may be refused on unsupported inputs. Availability,
construction success, empirical functionality, size, and diversity
are therefore distinct outcomes.

The benign x64 and x86 fixture results and the separate 102-payload
campaign appear in Section~\ref{sec:extended_results}. The latter
includes unbuilt and non-working outputs. Functional results on the
curated fixtures are not used as evidence of universal payload support.
The older selected x64 wrapping/static-scan cohort is also identified
separately, including its construction failures and duplicate outputs.
